\section{Lean 4 formalization outline and completion path}

The formalization should be layered.  This keeps ordinary algebraic geometry
separate from the final arithmetization and permits the field-level statements
to be reused independently.

\subsection{Layer A: generic field geometry}

Create the following modules under a new library, for example
\code{ConstructiveGeometryIPRA}.
\begin{description}
\item[\code{Vec3.lean}] triples, dot product, cross product, determinant,
  scalar proportionality, and the join--meet identities;
\item[\code{Projective.lean}] coded projective points and lines with explicit
  nonzero fields, incidence, collinearity, concurrency, and transport under
  invertible matrices;
\item[\code{Affine.lean}] the embedding $(x,y)\mapsto(x,y,1)$, affine
  combinations, signed brackets, line parameters, and directed ratios;
\item[\code{Conic.lean}] homogeneous quadratics, restrictions to coordinate
  sides, tangent covectors, the dual conic, and rational normalization;
\item[\code{PolynomialSide.lean}] univariate coefficient lists, evaluation,
  products of linear factors, and the certified-side-product lemma.
\end{description}

A representative kernel definition is:
\begin{lstlisting}[language={},caption={Suggested field-level kernel}]
structure V3 (K : Type*) where
  x : K
  y : K
  z : K

def cross [CommRing K] (u v : V3 K) : V3 K :=
  V3.mk (u.y*v.z-u.z*v.y)
        (u.z*v.x-u.x*v.z)
        (u.x*v.y-u.y*v.x)

def dot [CommRing K] (u v : V3 K) : K :=
  u.x*v.x + u.y*v.y + u.z*v.z

def det [CommRing K] (u v w : V3 K) : K := dot u (cross v w)
\end{lstlisting}
The core identities close with \code{ring}; nonzero and cancellation steps
remain explicit hypotheses.

\subsection{Layer B: one theorem module per certificate}

The suggested modules and decisive lemmas are:
\begin{longtable}{p{0.27\textwidth}p{0.63\textwidth}}
\toprule
Module & Main theorem or certificate \\
\midrule
\endhead
\code{PappusHarmonic.lean} & Compute $C,X,W$ after affine normalization and prove the bracket harmonic equation. \\
\code{CevaMenelaus.lean} & Expand the two determinants; for Ceva verify the explicit point, and for Menelaus verify the explicit line. \\
\code{CarnotConic.lean} & Prove the binary-quadratic lemma and the three explicit factorization identities for the converse witness. \\
\code{CarnotPolygon.lean} & Prove endpoint evaluation for a list product and telescope by induction over the polygon list. \\
\code{Hoehn.lean} & Derive the four-point determinant identity, the three cyclic relabellings, and the two five-factor cancellations. \\
\code{PappusHexagon.lean} & Prove the coordinatewise identity $X-Y+Z=0$. \\
\code{PascalBrianchon.lean} & Prove Pascal's weighted vector identity, dualize it syntactically, and verify the explicit concurrency point. \\
\bottomrule
\end{longtable}

The archive includes \code{lean/FieldCertificates.lean}, a Mathlib-oriented
source file containing the most compact scalar identities with \code{ring}
proofs.  It is a starting point, not claimed here as a build-tested module.

\subsection{Layer C: primitive-recursive rational codes}

Define code-level structures without quotient types:
\begin{lstlisting}[language={}]
structure RatCode where
  num : IntCode
  den : Nat
  den_pos : den > 0

def RatEq (a b : RatCode) : Prop :=
  intMul a.num (ofNat b.den) = intMul b.num (ofNat a.den)
\end{lstlisting}
Implement code operations and prove congruence with respect to
\code{RatEq}.  Then define coded triples, matrices, coefficient lists, cross
products, determinants, and finite products.  Each operation must be connected
to an explicit primitive-recursive symbol of the existing \code{PRAAx}
language.  A sparse-polynomial datatype with normalization is useful, but not
strictly required: because every fixed identity is small, direct equational
derivations are also viable.

\subsection{Layer D: syntactic \texorpdfstring{$\ipra$}{iPRA} proofs}

For each theorem:
\begin{enumerate}
\item write the exact \code{Fm PTm n} formula, including all well-formedness
and apartness premises;
\item construct the witness terms before proving the implication;
\item compile the field-level polynomial certificate into a derivation under
\code{PRAAx};
\item use quantifier-free induction only for list products and bounded-search
correctness; and
\item close existential conclusions by introduction with the explicit term.
\end{enumerate}

The final bridge already has the required shape:
\begin{lstlisting}[language={}]
import ConciseIPRA.BiInterpretation

open LeanProofs.ConciseIPRA

-- Given h : [] |-[PRAAx] phi,
-- (pra_theorem_iff phi).mp h has type
-- [] |-[TeqAx] phi.map (fun _ => FromPRA.sigma).
\end{lstlisting}
The archive's \code{lean/IPRAFormalizationPlan.lean} records the intended
module interfaces next to this import.

\subsection{Recommended implementation order}

The shortest dependency path is:
\begin{enumerate}
\item \code{Vec3}, then Pappus hexagon, Ceva, and Menelaus;
\item Pappus harmonic and Hoehn;
\item binary quadratics and Carnot conic;
\item polynomial lists and Carnot polygon;
\item Pascal on the standard conic, then dual Brianchon;
\item rational code correctness and the fixed-identity compiler;
\item theorem-by-theorem \code{PRAAx} derivations; and
\item transport to \code{TeqAx} and, optionally, to a smaller native calculus.
\end{enumerate}
This order gives useful completed milestones early and postpones only the
unavoidably variable list infrastructure.

\subsection{Completion criteria}

A complete formalization must satisfy all of the following:
\begin{itemize}
\item no \code{sorry}, \code{admit}, unproved axiom, or semantic appeal to a
field tactic inside the final arithmetic derivation;
\item successful build with the repository's pinned Lean and Mathlib versions;
\item explicit formulas for all eight arithmetized statements;
\item kernel-checked \code{PRAAx} derivations, including Carnot's list
inductions and the bounded searches used by normalization;
\item an application of the existing bi-interpretation theorem to each final
statement; and
\item acceptance by the chosen native proof checker if a compressed
certificate is exported.
\end{itemize}
The present package reaches the human-proof and executable-algebra stages and
specifies this remaining path without claiming that those final criteria have
already been met.
